\documentclass[10pt,a4paper]{amsart}
\usepackage[margin=19mm]{geometry}
\usepackage{amsmath,amssymb,mathtools}
\usepackage[scr=rsfs,cal=euler]{mathalfa}
\usepackage{xfrac}
\usepackage[unicode,hidelinks,bookmarks=false]{hyperref}
\newcommand{\ie}{i.e.\ }
\newcommand{\cf}{cf.\ }
\newcommand{\resp}{respectively\ }
\newcommand{\Subsection}{\subsection}
\providecommand{\nobreakdash}{\nobreak\hbox{-}}
\setlength{\emergencystretch}{2em}
\multlinegap0pt
\usepackage{amssymb}
\usepackage{enumerate}
\usepackage{tikz}
\usepackage[normalem]{ulem}
\usepackage{algorithm}\usepackage{algpseudocode}
\usepackage{bm}
\usepackage{float}
\newtheorem{prop}{Proposition}
\DeclareMathOperator*{\argmax}{\ensuremath{arg\,max}}
\DeclareMathOperator*{\argmin}{\ensuremath{arg\,min}}
\DeclareMathOperator*{\argopt}{\ensuremath{arg\,opt}}
\newcommand{\cut}{\mathrm{cut}}
\newcommand{\labeling}{\ensuremath{\mathrm{lab}}}
\DeclareMathOperator*{\opt}{\ensuremath{opt}}
\title{A Note on $1$-spectral clustering}
\author{Sihong Shao, Chuan Yang}
\date{Source: arXiv:2607.17158v1 (CC BY 4.0)}
\begin{document}
\maketitle

\noindent\emph{Source and licence.} A Note on $1$-spectral clustering. Version arXiv:2607.17158v1 (CC BY 4.0), distributed by arXiv under the Creative Commons Attribution 4.0 International licence, http://creativecommons.org/licenses/by/4.0/, as recorded in the arXiv licence metadata for this exact version and shown on the abstract page https://arxiv.org/abs/2607.17158v1. Full licence notice: \texttt{licenses/arxiv-2607.17158-CC-BY-4.0.txt}.\par\smallskip

\noindent\textit{Editorial scope.} Counted unit: the article's prop thm::maxkcut (source0.tex lines 1673--1676). The article's complete original proof of it is delivered with it (source0.tex lines 1677--1680, one \texttt{proof} environment of 4 lines). No mathematics is written, repaired or re-derived: the statement and the proof are the article's own lines, in the article's own notation, and no retained line is substituted or re-flowed.  Retained uncounted as the source-local context the proof needs: lines 1506--1531.  Nothing was omitted from any retained span, so the package carries no omission record.  No retained line contains a citation, so the wrapper carries no bibliography and no bibliography entry.  No retained line contains a figure environment, an image-inclusion command or drawing code, so no image is delivered and the package declares no asset dependency.  Editorial formatting only, in addition: an AMS \texttt{amsart} preamble that restates the article's own declarations and macros used by the retained lines, namely the environments \texttt{prop} on separate counters, together with the standard packages those lines need.  The retained environments print in this document as Proposition 1 in order of appearance, whereas the article prints the same items with the numbering of its own full document.  The complete native source of the article as distributed in the arXiv e-print for this exact version is bundled unchanged as \texttt{sources/205539/source0.tex}.
			\textbf{The single-transfer local search operator $\widetilde{O}_1$ or $\widehat{O}_1$}  moves the vertex ${v}^*$ to $S_{{q}^*}$ that satisfies
			%will select the vertex $\widetilde{v}^*$ and transfer it into the subset $S_{\widetilde{q}^*}$ using the following procedure:
			\begin{align}
				\label{choice::o1cut}
				({v}^*,\,{q}^*)
				&\in 
				\argopt_{(v,q)\in \mathbb{D}(I)}\left\{F(I\circ (v\rightarrow S_q))\right\},
				%					\left\{
				%					\begin{aligned}
					%						&\argopt_{(v,q)\in \mathbb{D}(I)}\left\{F(I\circ (v\rightarrow S_q))\right\}, &
					%						\text{if }\opt\, F\in\text{MaxGCP},\\
					%						&\argopt_{(v,q)\in\mathbb{D}(I)}\left\{F(I\circ (v\rightarrow S_q))\right\}, &
					%						\text{if }\opt\, F\in\text{MinGCP},
					%					\end{aligned}
				%					\right. 
				\\
				\label{eq::d_vq}
				\mathbb{D}(I)
				&=\left\{
				\begin{aligned}
					\bigcup\limits_{t\in [k]\backslash\{\labeling(v, I)\}}\argmax_{v\in V,q=t}\{\mathcal{D}_{vq}\}, & \text{~if~}\opt\, F\in\text{MaxGCP},\\
					\bigcup\limits_{t\in [k]\backslash\{\labeling(v, I)\}}\argmin_{v\in V,q=t}\{\mathcal{D}_{vq}\},  & \text{~if~}\opt\, F\in\text{MinGCP}.
				\end{aligned}
				\right.
			\end{align}
			It can be readily verified that the search area $\mathbb{D}(I)$ contains the simplest moves leading to the largest improvement in the cut value or boundary size, thereby indicating that $\mathbb{D}(I)$ keeps pace with the basic needs of $\opt\,F$ and is thus the promising search area of improved solutions for $F$. In a word, $\widetilde{O}_1$ or $\widehat{O}_1$ selects the best simplest move within $\mathbb{D}(I)$, a subset of all possible moves in the first order neighborhood of the current solution $I$.

			\begin{prop}
				\label{thm::maxkcut}
				Given any unweighted graph and ``$F=\frac{1}{2}\sum_{p=1}^k\cut(\partial S_p)$", the operators $\widetilde{O}_1$ and $\widetilde{O}_2$ are able to identify local optima in the first and second order neighborhoods, respectively.
			\end{prop}

			\begin{proof}
				\textsc{Min}-$k$-\textsc{Cut} shares the similar proof with \textsc{Max}-$k$-\textsc{Cut} that we assume ``$\opt=\max$" in the followings.
				If $\widetilde{O}_1$ can not improve the cut value for the current solution $I$, $\widetilde{D}_{vq}\leq 0$ holds for $\forall\,(v,q)\in \mathbb{D}(I)$ (see Eq.~\eqref{eq::d_vq}), thus $\widetilde{D}_{vq}\leq 0$ holds for $\forall\,v\in V$ and $\forall\,q\in[k]$, leading to $I$ being the local optimum in the first order neighborhood. Next we prove that if there exists a double-transfer move $S_{p_u}\rightarrow u\rightarrow S_{q_u}$ and $S_{p_v}\rightarrow v \rightarrow S_{q_v}$ that can improve the cut value, i.e. $\widetilde{\Delta}_{(u,v)}>0$, $\widetilde{O}_2$ will promote the cut value. Since the current solution $I$ is a local maximum in $\mathbb{E}^{(1)}(I)$, we have $\widetilde{\Delta}_{u\rightarrow S_{q_u}}\leq 0$ and $\widetilde{\Delta}_{v\rightarrow S_{q_v}}\leq 0$, indicating $\widetilde{\phi}w_{uv}=2$ to satisfy $\widetilde{\Delta}_{(u,v)}>0$ and at least one of $\widetilde{\Delta}_{u\rightarrow S_{q_u}}=0$ and $\widetilde{\Delta}_{v\rightarrow S_{q_v}}=0$ must be satisfied, implying that either $(v,q_v)$ or $(u,q_u)$ belongs to $\mathbb{D}(I)$. Therefore, this double-transfer move $S_{p_u}\rightarrow u\rightarrow S_{q_u}$ and $S_{p_v}\rightarrow v \rightarrow S_{q_v}$ is within the search area of $\widetilde{O}_2$.
			\end{proof}

\end{document}
